\documentclass[11pt]{article}

\usepackage[margin=1in]{geometry}
\usepackage{csquotes}

\title{Demographic transformation, liberal democracy, and the crisis of political judgment}
\author{Carlo Perassi}
\date{June 8, 2026}

\begin{document}

\maketitle

\begin{abstract}
This essay offers a polemical critique of contemporary Western immigration policy, with particular attention to the perceived failure of political elites to distinguish between ordinary migration, successful integration, and the importation of illiberal or theocratic political cultures. The argument is not that migrants, Muslims, or any religious community should be treated as collectively guilty. Rather, it is that liberal democracies may endanger themselves when they refuse to defend the cultural, legal, and institutional preconditions of their own survival.
\end{abstract}

%\tableofcontents

\section{Introduction}

A recurring anxiety in contemporary Western politics concerns the relation between demographic change, cultural integration, and the preservation of liberal democracy. In Europe and North America, large-scale migration has often been justified through humanitarian, economic, and demographic arguments. Yet these justifications become fragile when political institutions are unwilling to address a central question: can a liberal democratic order survive if it admits, in significant numbers, populations or ideological movements that do not accept the moral and legal premises of that order?

The issue is not migration as such. Nor is it religion as such. The issue is whether a society can remain liberal while tolerating the growth of explicitly illiberal, theocratic, or anti-democratic norms within its own borders.

\section{The alleged failure of the Western left}

A severe criticism may be directed at sections of the Western left. In the name of anti-racism, anti-colonialism, humanitarianism, or hostility toward what they call ``the West'', they appear, at times, to have lost the ability to defend the very liberal institutions that allow them to speak, organize, and criticize freely.

This produces a striking contradiction. The same political culture that often invoked Karl Popper's principle that a tolerant society must not be infinitely tolerant toward the intolerant seemed willing to apply that principle forcefully against neo-fascist or neo-nazi movements. Yet it often hesitates to apply the same principle to religious or political movements that reject freedom of conscience, equality before the law, freedom of expression, women's rights, secular law and democratic pluralism.

The question is therefore unavoidable: why should one form of anti-liberalism be treated as an existential threat, while another is excused, minimized, or explained away?

\section{Integration and the limits of demographic substitution}

Some defenders of mass immigration argue that demographic decline in Europe and North America requires large-scale population replacement through immigration. This argument is often presented as economically pragmatic: ageing societies need workers, taxpayers and demographic renewal.

However, demographic substitution is not a mechanical process. A society is not sustained merely by replacing one population with another. It also depends on inherited norms of trust, work, legality, reciprocity, civic responsibility and institutional loyalty. These norms cannot be assumed to reproduce themselves automatically.

For example, if a region has developed a strong work culture, a high-trust civil society and relatively efficient institutions, these features are not guaranteed to survive under conditions of rapid demographic and cultural transformation. They require transmission, education, assimilation and a shared acceptance of the common polity rules.

A society that treats population numbers as interchangeable while neglecting cultural integration risks undermining the very foundations of its economic and political life.

\section{The danger of parallel societies}

A liberal democracy may tolerate pluralism but it cannot indefinitely tolerate the formation of parallel legal and moral orders that reject the authority of democratic law. Voluntary segregation, religiously or ethnically defined enclaves, and informal communities governed by anti-liberal norms may begin as local phenomena. Over time, however, they may expand into durable parallel societies.

The danger is not merely cultural discomfort. It is institutional fragmentation. When democratic law becomes only one normative system among others, and when certain communities increasingly defer to religious, clan-based, or extra-legal authorities, the state loses its monopoly on legitimate civic authority.

In such a scenario, political actors may eventually argue that limited concessions to religious law, including elements inspired by Sharia, are reasonable, pragmatic, or necessary for social peace. What begins as tolerance may gradually become capitulation.

\section{The paradox of liberal self-destruction}

History contains many examples of civilizations destroyed by conquest, collapse, disease, invasion, or internal decay. What appears more unusual is the possibility of a civilization actively facilitating its own dissolution while describing this process as moral progress.

This phenomenon may be interpreted as a form of political self-negation. A society inherits freedom, secular legality, scientific reason, artistic creativity, and individual rights, yet becomes so ashamed of itself that it no longer wishes to defend the conditions that produced those achievements.

This is especially paradoxical when the forces being tolerated are not more liberal, more rational, or more humane than the society being criticized, but in many respects less so. The result is not liberation, but the risk of regression toward social forms characterized by coercion, censorship, patriarchy, religious domination, and the suppression of dissent.

\section{A secular objection}

This critique need not be Christian. One may reject Christianity, or stand outside religious belief entirely, and still recognize that the liberal West emerged from a long and complex historical struggle involving Christianity, secularization, Enlightenment rationalism, constitutionalism, and scientific culture.

Bertrand Russell, for example, was not a Christian believer, yet he understood the importance of reason, freedom of thought, and moral criticism. A secular person may therefore oppose theocratic politics not because of loyalty to Christianity, but because of loyalty to liberty, reason, and human dignity.

The question is not whether one religion should defeat another. The question is whether secular liberal civilization has the will to defend itself against every ideology that seeks to subordinate individual freedom to dogma.

\section{The moral horror of gradual decline}

The most disturbing aspect of this process is not that collapse may happen suddenly. It is that it may happen gradually, one concession at a time. A society may continue to speak the language of rights while emptying those rights of substance. It may continue to praise democracy while tolerating communities and ideologies that despise democracy. It may continue to celebrate freedom while teaching its citizens that naming obvious problems is itself forbidden.

In such a context, public life becomes marked by denial. Many people may see the contradiction clearly, yet fear saying so. The emperor is naked, but the social cost of stating the obvious becomes too high.

\section{Conclusion}

A liberal democracy cannot survive by hatred, paranoia, or collective blame. But neither can it survive through cowardice, denial, or moral confusion. It must distinguish between persons and ideologies, between peaceful religious belief and political theocracy, between immigration and failed integration, between pluralism and self-dissolution.

The central claim is therefore not that outsiders are inherently enemies, nor that any group is collectively incapable of integration. The claim is that Western societies must recover the confidence to defend liberal democracy as a substantive civilizational achievement. If they refuse to do so, they may discover too late that tolerance without judgment becomes the instrument of its own destruction.

\end{document}
